\section*{अध्याय \ref{ch:g12:photosynthesis} --- प्रकाशसंश्लेषण}

\begin{solution}{exo:g12:photosynthesis:1}
दोहरी झिल्ली वाला एक \omterm{def:g10:cells-common-unit:organelle}{कोशिकांग}, जिसमें तरल पीठिका है और उसमें \omterm{def:g12:photosynthesis:chloroplast}{थाइलेकॉइड}
झिल्लियों के चट्टे हैं। वर्णक और प्रकाश बदलते \omterm{def:g11:gene-expression:protein}{प्रोटीन} \omterm{def:g12:photosynthesis:chloroplast}{थाइलेकॉइड} झिल्लियों
में हैं; और शर्करा बनाते \omterm{def:g11:enzymes-and-phenotype:enzyme}{एंजाइम} पीठिका में।
\end{solution}

\begin{solution}{exo:g12:photosynthesis:2}
पर्णहरित नीला तथा लाल अवशोषित करता है और हरा परावर्तित करता है, और वही हमें
दिखाई देता है। नीला तथा लाल प्रकाश \omterm{def:g10:cell-metabolism:autotrophy}{प्रकाशसंश्लेषण} सबसे अच्छा चलाते हैं।
\end{solution}

\begin{solution}{exo:g12:photosynthesis:3}
ATP, अपचयित NADP और ऑक्सीजन; ऑक्सीजन उपोत्पाद है, जो छोड़ दी जाती है।
\end{solution}

\begin{solution}{exo:g12:photosynthesis:4}
चीरे गए पानी से, कार्बन डाइऑक्साइड से नहीं: भारी ऑक्सीजन से चिह्नित पानी
पाए शैवालों ने भारी $\mathrm{O_2}$ छोड़ी; और चिह्नित कार्बन डाइऑक्साइड पाने
पर नहीं छोड़ी।
\end{solution}

\begin{solution}{exo:g12:photosynthesis:5}
छह $\mathrm{CO_2}$, 18 ATP और 12 अपचयित NADP।
\end{solution}

\begin{solution}{exo:g12:photosynthesis:6}
\qty{550}{nm} पर: अवशोषण लगभग 8\%, दर लगभग 25\%। और \qty{665}{nm} पर:
अवशोषण लगभग 90\%, दर लगभग 95\%। हरे लैंप वही प्रकाश पहुँचाते हैं जिसे वह
पत्ता सबसे कम इस्तेमाल करता है; वह हरितगृह लगभग कुछ भी न उगाकर ऊर्जा बचाता।
\end{solution}

\begin{solution}{exo:g12:photosynthesis:7}
वे जीवाणु ऑक्सीजन की ओर तैरते हैं; वे वहीं जमा होते हैं जहाँ ऑक्सीजन छोड़ी
जा रही हो, यानी जहाँ वह शैवाल \omterm{def:g10:cell-metabolism:autotrophy}{प्रकाशसंश्लेषण} करता हो; और वह स्पेक्ट्रम के
नीले तथा लाल हिस्सों के नीचे है। \omterm{def:g10:cell-metabolism:autotrophy}{प्रकाशसंश्लेषण}, जो अपनी ऑक्सीजन से नापा
गया, प्रकाश के रंग पर निर्भर है और उन वर्णकों के अवशोषण के पीछे चलता है।
\end{solution}

\begin{solution}{exo:g12:photosynthesis:8}
बिना किसी कार्बन के बँधे ऑक्सीजन छोड़ी जाती है: यानी प्रकाश प्रावस्था, जो
पानी चीरती और इलेक्ट्रॉन हिलाती है, कार्बन प्रावस्था से अलग है और उसे
$\mathrm{CO_2}$ नहीं चाहिए; उसे प्रकाश, पानी और इलेक्ट्रॉन लेने को कोई चीज़
चाहिए।
\end{solution}

\begin{solution}{exo:g12:photosynthesis:9}
प्रकाश के बिना न कोई ATP है न कोई अपचयित NADP, इसलिए वह तीन-कार्बन अम्ल
अपचयित नहीं हो पाता और जमा होता है; और वह ग्राही कुछ देर $\mathrm{CO_2}$
पकड़ता रहता है पर दोबारा नहीं बनता, तथा चुक जाता है।
\end{solution}

\begin{solution}{exo:g12:photosynthesis:10}
प्रकाश तीव्रता दूरी के साथ गिरती है (मोटे तौर पर उसके वर्ग के हिसाब से: दुगनी
दूरी पर चार गुना कम), और प्रकाश प्रावस्था धीमी चलती है। अँधेरे में कोई
ऑक्सीजन नहीं छोड़ी जाती --- और वह पौधा श्वसन से कुछ खपा भी लेता है।
\end{solution}

\begin{solution}{exo:g12:photosynthesis:11}
$8/44 \approx \qty{0.18}{mol}$ $\mathrm{CO_2}$: यानी $0.18/6 \approx
\qty{0.030}{mol}$ ग्लूकोज, लगभग \qty{5.5}{g}; और \qty{0.18}{mol}
$\mathrm{O_2}$, लगभग \qty{5.8}{g}।
\end{solution}

\begin{solution}{exo:g12:photosynthesis:12}
वह ख़ुद प्रकाश इस्तेमाल नहीं करती, केवल प्रकाश प्रावस्था के उत्पाद। पर
अँधेरे में रखे किसी पौधे में वह एक मिनट के भीतर रुक जाती है, जैसे ही वे ATP
तथा अपचयित NADP ख़र्च हो जाएँ: वह प्रकाश में चलती है, प्रकाश प्रावस्था से
खिलाई हुई, और ``अँधेरी'' यह बताता है कि उसे क्या चाहिए, यह नहीं कि वह कब
होती है।
\end{solution}

\begin{solution}{exo:g12:photosynthesis:13}
$^{18}$O: एक मिनट के भीतर छोड़ी गई $\mathrm{O_2}$ में, और एक घंटे तथा एक दिन
बाद भी वहीं (कुछ श्वसन से बने पानी में भी)। $^{14}$C: एक मिनट बाद उस चक्र के
तीन-कार्बन अम्ल तथा शर्कराओं में; एक घंटे बाद सुक्रोज तथा मंड में; और एक दिन
बाद सेल्युलोज़, \omterm{def:g10:chemistry-of-life:families}{लिपिड}, \omterm{def:g10:chemistry-of-life:families}{ऐमीनो अम्लों} में, तथा श्वसन से निकली $\mathrm{CO_2}$
में।
\end{solution}

\begin{solution}{exo:g12:photosynthesis:14}
$20/1000 = 2\%$। बाक़ी परावर्तित होता है (वह हरा), ग़लत तरंगदैर्घ्यों में
अवशोषित होकर ऊष्मा बन जाता है, उस वाष्पोत्सर्जन में ख़र्च होता है जो पानी
वाष्पित करता है, या ख़ुद प्रकाश अभिक्रियाओं में ऊष्मा की तरह खो जाता है।
\end{solution}

\begin{solution}{exo:g12:photosynthesis:15}
श्वसन, दहन तथा चट्टानों और लोहे का ऑक्सीकरण लगातार ऑक्सीजन खपाते हैं; और
उसे नया करते \omterm{def:g10:cell-metabolism:autotrophy}{प्रकाशसंश्लेषण} के बिना वायुमंडल का भंडार कुछ ही हज़ार सालों में
चुक जाता। शुरुआती वायुमंडल में कोई मुक्त ऑक्सीजन नहीं थी; वह तभी जमा हुई जब
प्रकाशसंश्लेषी जीव उसे बहुत लंबे समय तक बना चुके थे --- यानी पहले दो अरब
साल ऑक्सीजन में ग़रीब थे।
\end{solution}

\begin{solution}{pb:g12:photosynthesis:1}
\textbf{1.} $8/44 \approx \qty{0.18}{mol}$।

\textbf{2.} $0.18/6 = \qty{0.030}{mol}$, लगभग \qty{5.5}{g}।

\textbf{3.} \qty{0.18}{mol} $\mathrm{O_2}$, लगभग \qty{4.4}{L}।

\textbf{4.} प्रति $\mathrm{O_2}$ पानी के दो अणु: \qty{0.36}{mol},
लगभग \qty{6.5}{g} --- यानी वाष्पोत्सर्जित \qty{600}{mL} का एक प्रतिशत।

\textbf{5.} \omterm{def:g12:photosynthesis:chloroplast}{थाइलेकॉइडों} में चीरे गए पानी के अणुओं से: जो भारी-ऑक्सीजन
चिह्नन से दिखाया गया, क्योंकि वह चिह्न $\mathrm{O_2}$ में तभी प्रकट होता है
जब पानी चिह्नित हो।

\textbf{6.} प्रति $\mathrm{CO_2}$ एक फेरा: \qty{0.18}{mol} फेरे।

\textbf{7.} प्रति फेरा तीन ATP और दो अपचयित NADP: \qty{0.55}{mol}
ATP और \qty{0.36}{mol} अपचयित NADP।

\textbf{8.} हर चीरा पानी दो इलेक्ट्रॉन देता है, यानी एक अपचयित NADP:
\qty{0.36}{mol} पानी --- यानी सवाल 4 वाला ही आँकड़ा, जैसा होना ही
चाहिए।

\textbf{9.} ऑक्सीजन तुरंत रुक जाती; अपचयित NADP तथा ATP बनना बंद हो जाते और
सेकंडों के भीतर चुक जाते; और शर्करा उसी वक़्त रुक जाती जब वे वाहक ख़त्म हो
जाएँ।

\textbf{10.} वे वाहक अब \omterm{prop:g12:photosynthesis:calvin}{कैल्विन चक्र} से ख़र्च नहीं होते; और उनके इलेक्ट्रॉन
तथा ऊर्जा लेने को कुछ न होने पर प्रकाश प्रावस्था सेकंडों के भीतर जाम होकर
धीमी पड़ जाती --- यानी दोनों प्रावस्थाएँ उन वाहकों से जुड़ी हैं।

\textbf{11.} प्रति घंटा $0.030 \times 2870 \approx \qty{86}{kJ}$, यानी
$86000/3600 \approx \qty{24}{W}$।

\textbf{12.} पूरी धूप का $24/1000 = 2.4\%$; और अवशोषित तरंगदैर्घ्यों का
$24/450 \approx
5\%$।

\textbf{13.} फ़ोटॉनों के $8 \times 176 \approx \qty{1400}{kJ}$ प्रति मोल
$\mathrm{CO_2}$; और जमा: $2870/6 \approx \qty{480}{kJ}$।

\textbf{14.} लगभग एक तिहाई; बाक़ी दोनों प्रावस्थाओं के एक के बाद एक ऊर्जा
अंतरणों में ऊष्मा की तरह छूट जाता है।

\textbf{15.} शुद्ध लाभ प्रति घंटा $5.5 - 0.5 = \qty{5}{g}$ ग्लूकोज।
प्रकाश में शुद्ध $\mathrm{CO_2}$ ग्रहण की तरह नापा जाता है; और अकेला
श्वसन अँधेरे में नापकर वापस जोड़ा जाता है ताकि सकल \omterm{def:g10:cell-metabolism:autotrophy}{प्रकाशसंश्लेषण} मिले।

\textbf{16.} $\num{8.5e11}/\num{1e11} \approx 8.5$ साल --- और समुद्र के
शैवाल उतना ही और बाँधें तो लगभग 4 साल।

\textbf{17.} श्वसन (पौधों, जंतुओं, कवकों तथा जीवाणुओं का), आग के साथ,
कार्बन डाइऑक्साइड लगभग उसी दर पर लौटाता है; बंधन और वापसी संतुलित थे, इसलिए
वह भंडार तब तक स्थिर था जब तक जीवाश्म ईंधन के जलने ने एक नया स्रोत नहीं जोड़
दिया।

\textbf{18.} \qty{1e11}{t} कार्बन $\num{8.3e12}\,\unit{mol}$ है; और उतने ही
मोल $\mathrm{O_2}$ लगभग \qty{2.7e11}{t} ऑक्सीजन हैं।

\textbf{19.} उत्पादन के $\num{1.2e15}/\num{2.7e11} \approx 4500$ साल
(समुद्र के हिस्से समेत लगभग 2000)। पर श्वसन ऑक्सीजन लगभग उसी दर पर खपाता है
जिस दर पर \omterm{def:g10:cell-metabolism:autotrophy}{प्रकाशसंश्लेषण} उसे बनाता है, इसलिए वह भंडार एक संतुलन है, कोई
जमाव नहीं: हवा की ऑक्सीजन उत्पादन के खपत से थोड़े-से आधिक्य से करोड़ों सालों
पर बनी।

\textbf{20.} पूरी धूप का लगभग 2\% (काम आती तरंगदैर्घ्यों का 5\%) शर्करा की
तरह जमा होता है; छोड़ी गई ऑक्सीजन चीरे गए पानी की ऑक्सीजन है; और वायुमंडल
के कार्बन का मोटे तौर पर एक चौथाई हर साल \omterm{def:g10:cell-metabolism:autotrophy}{प्रकाशसंश्लेषण} से होकर गुज़रता है।
\end{solution}
